\section{The canonical multiwave Hamiltonian}
\label{sec:canonical}

We first give the complex canonical system in the three-matrix notation
of \cref{eq:intro-mixing}. An explicit change of coordinates identifies
its cubic interactions with the multiwave construction of Jiang and
Rauch-Wojciechowski \cite[Section II]{JRW1991}. Retaining independent
constant diagonal blocks supplies the mode-dependent linear rates
required by compatible prescribed BA data in
\cref{thm:general-BA-realization}. We refer to this constant-diagonal
extension as the modified JRW system. Commuting power charges then redistribute
these rates into the symmetric phase convention used here. The resulting
pencil supplies commuting integrals through the AKS construction;
their functional independence is considered separately. The detailed
comparison with the original system is given in Appendix~\ref{sec:JRW-dictionary}.

\subsection{Modes, brackets and Hamilton equations}
\label{sec:canonical-modes}

There are $p=LM+MN+NL$ canonical mode pairs. Their nonzero basic
brackets are
\begin{equation}\label{eq:canonical-brackets}
\begin{aligned}
 \{u_{\ell m},\bar{u}_{ij}\}&=-i\delta_{\ell i}\delta_{mj},
 &&1\leq\ell,i\leq L,\quad1\leq m,j\leq M,\\
 \{v_{mn},\bar{v}_{jk}\}&=-i\delta_{mj}\delta_{nk},
 &&1\leq m,j\leq M,\quad1\leq n,k\leq N,\\
 \{w_{n\ell},\bar{w}_{ki}\}&=-i\delta_{nk}\delta_{\ell i},
 &&1\leq n,k\leq N,\quad1\leq\ell,i\leq L.
\end{aligned}
\end{equation}
All other basic brackets vanish, apart from antisymmetry. Bars label
independent Hamiltonian partners, with the same array dimensions as the
unbarred modes. Stars mean actual complex conjugation; $\trans$ means
ordinary transpose, and $\dagger$ means actual conjugate transpose.
A dot denotes differentiation with respect to time $t$, and
$\dot h=\{h,H\}$. Thus
\[
 i\dot u=\frac{\partial H}{\partial\bar u},\qquad
 i\dot{\bar u}=-\frac{\partial H}{\partial u}.
\]
This is the usual complex canonical convention: writing
$u=(q+ip)/\sqrt2$ and $\bar u=(q-ip)/\sqrt2$ gives $\{q,p\}=1$.
The variables $q,p$ are complex unless a real form is imposed.

For real phase coefficients and real coupling, the real forms
$\bar u=\epsilon_u u^*$, $\bar v=\epsilon_v v^*$ and
$\bar w=\epsilon_w w^*$, with $\epsilon_u,\epsilon_v,\epsilon_w\in\{1,-1\}$,
are invariant precisely when $\epsilon_u\epsilon_v\epsilon_w=1$.
The all-positive choice is the explosive three-wave form; the three
choices with two negative signs are decay forms and admit definite
combinations of the power charges. The Hamiltonian is real on each
form, with $i\dot u=\epsilon_u\partial H/\partial u^*$ and the
corresponding equations for $v,w$. The general complex system retains
independent partners; Liouville integrability is understood in the
complex Hamiltonian sense.

Let $b_\ell^{L}$, $b_m^{M}$ and $b_n^{N}$ be constants, with
$1\leq\ell\leq L$, $1\leq m\leq M$, $1\leq n\leq N$.
The superscripts distinguish the three groups of phase coefficients.

\begin{samepage}
With coupling normalization $\kappa=\sqrt2$, the Hamiltonian is
\begin{align}\label{eq:canonical-H}
H={}&\sum_{\ell=1}^{L}\sum_{m=1}^{M}
 (b_\ell^{L}+b_m^{M})u_{\ell m}\bar{u}_{\ell m}
 +\sum_{m=1}^{M}\sum_{n=1}^{N}
 (b_m^{M}+b_n^{N})v_{mn}\bar{v}_{mn}\nonumber\\
&+\sum_{n=1}^{N}\sum_{\ell=1}^{L}
 (b_n^{N}+b_\ell^{L})w_{n\ell}\bar{w}_{n\ell}\nonumber\\
&+\sqrt2\sum_{\ell=1}^{L}\sum_{m=1}^{M}\sum_{n=1}^{N}
 \left(u_{\ell m}v_{mn}w_{n\ell}
       +\bar{u}_{\ell m}\bar{v}_{mn}\bar{w}_{n\ell}\right).
\end{align}
\end{samepage}
Each mode's linear phase coefficient is the sum of the constants at its two
endpoint groups. The Hamiltonian is invariant in form under cyclic
permutation of the groups and their mode matrices.
The three mode powers are
\begin{equation}\label{eq:canonical-powers}
 I^u_{\ell m}=u_{\ell m}\bar{u}_{\ell m},\qquad
 I^v_{mn}=v_{mn}\bar{v}_{mn},\qquad
 I^w_{n\ell}=w_{n\ell}\bar{w}_{n\ell},
 \quad1\leq\ell\leq L,\ 1\leq m\leq M,\ 1\leq n\leq N.
\end{equation}
They become squared moduli when the partners are actual complex
conjugates.

The Hamilton equations are, for
$1\leq\ell\leq L$, $1\leq m\leq M$, $1\leq n\leq N$,
\begin{equation}\label{eq:canonical-rates}
\begin{aligned}
 i\dot u_{\ell m}&=(b_\ell^L+b_m^M)u_{\ell m}
            +\sqrt2\sum_{k=1}^{N}\bar{v}_{mk}\bar{w}_{k\ell},\\
 i\dot{\bar{u}}_{\ell m}&=-(b_\ell^L+b_m^M)\bar{u}_{\ell m}
            -\sqrt2\sum_{k=1}^{N}v_{mk}w_{k\ell},\\
 i\dot v_{mn}&=(b_m^{M}+b_n^N)v_{mn}
       +\sqrt2\sum_{i=1}^{L}\bar{u}_{im}\bar{w}_{ni},\\
 i\dot{\bar{v}}_{mn}&=-(b_m^{M}+b_n^N)\bar{v}_{mn}
       -\sqrt2\sum_{i=1}^{L}u_{im}w_{ni},\\
 i\dot w_{n\ell}&=(b_n^N+b_\ell^{L})w_{n\ell}
       +\sqrt2\sum_{j=1}^{M}\bar{u}_{\ell j}\bar{v}_{jn},\\
 i\dot{\bar{w}}_{n\ell}&=-(b_n^N+b_\ell^{L})\bar{w}_{n\ell}
       -\sqrt2\sum_{j=1}^{M}u_{\ell j}v_{jn}.
\end{aligned}
\end{equation}

\subsection{Power symmetries and phase conventions}
\label{sec:canonical-phase-gauges}

The linear terms cancel from the mode-power derivatives. The conserved
power combinations generating diagonal mode gauges are
\begin{equation}\label{eq:canonical-charges}
\begin{aligned}
 Q_\ell^L&=\sum_{n=1}^{N}I^w_{n\ell}
             -\sum_{m=1}^{M}I^u_{\ell m},&&1\leq\ell\leq L,\\
 Q_m^M&=\sum_{n=1}^{N}I^v_{mn}
             -\sum_{\ell=1}^{L}I^u_{\ell m},&&1\leq m\leq M,\\
 Q_n^N&=\sum_{\ell=1}^{L}I^w_{n\ell}
             -\sum_{m=1}^{M}I^v_{mn},&&1\leq n\leq N.
\end{aligned}
\end{equation}
They commute, and their only linear relation is
\begin{equation}\label{eq:canonical-charge-relation}
 \sum_{\ell=1}^{L}Q_\ell^L-\sum_{m=1}^{M}Q_m^M
                         -\sum_{n=1}^{N}Q_n^N=0.
\end{equation}
Hence there are $L+M+N-1$ independent power charges on a generic
mode chart. The mixed signs in their definitions specify the
orientation of each gauge action.

Adding a constant linear combination of these charges to $H$
redistributes the linear mode rates without changing any power
derivative. We distinguish the symmetric-phase Hamiltonian $H$ from
a second convention $H_0$, in which the $u$ modes have zero linear
rate. This convention is convenient for comparison with Jiang and
Rauch-Wojciechowski. Define
\begin{align}\label{eq:canonical-H0}
 H_0&=H+\sum_{\ell=1}^{L}b_\ell^LQ_\ell^L
        +\sum_{m=1}^{M}b_m^MQ_m^M\nonumber\\
 &=\sum_{n=1}^{N}\sum_{\ell=1}^{L}
       (2b_\ell^L+b_n^N)w_{n\ell}\bar{w}_{n\ell}
   +\sum_{m=1}^{M}\sum_{n=1}^{N}
       (2b_m^M+b_n^N)v_{mn}\bar{v}_{mn}+H_{\mathrm{mix}},
\end{align}
where $H_{\mathrm{mix}}$ is \cref{eq:intro-mixing} with
$\kappa=\sqrt2$. Let
$(\tilde{\boldsymbol u},\tilde{\boldsymbol v},\tilde{\boldsymbol w})$
and their barred partners denote solutions of Hamilton's equations
generated by $H_0$. The corresponding solutions of $H$ are
\begin{equation}\label{eq:canonical-phase-gauge}
\begin{aligned}
 u_{\ell m}(t)&=e^{-i(b_\ell^L+b_m^M)t}\tilde u_{\ell m}(t),
 &&1\leq\ell\leq L,\quad1\leq m\leq M,\\
 v_{mn}(t)&=e^{ib_m^Mt}\tilde v_{mn}(t),
 &&1\leq m\leq M,\quad1\leq n\leq N,\\
 w_{n\ell}(t)&=e^{ib_\ell^Lt}\tilde w_{n\ell}(t),
 &&1\leq n\leq N,\quad1\leq\ell\leq L.
\end{aligned}
\end{equation}
The barred partners receive the
inverse factors. This time-dependent change is canonical at each
fixed time. Its exponents cancel in every cubic mixing product, so
both those products and all mode powers are identical in the two
descriptions. Appendix~\ref{sec:JRW-dictionary} identifies $H_0$ with the
constant-diagonal extension of the 1991 Hamiltonian.

The independent diagonal parameters and this change of phase convention
have distinct roles. In each mixing product $u_{\ell m}v_{mn}w_{n\ell}$,
the sum of the three linear phase coefficients is
$2(b_\ell^L+b_m^M+b_n^N)$, for
$1\leq\ell\leq L$, $1\leq m\leq M$, $1\leq n\leq N$.
The phase gauge \cref{eq:canonical-phase-gauge} preserves this sum.
The outer diagonal parameters therefore supply interaction-dependent
linear rates that remain present in either convention. Their values
are fixed by the marked curve data in the BA realization.

\subsection{The matrix pencil and its Lax equation}
\label{sec:canonical-lax}

A Lax representation packages the mode equations into a matrix
commutator, from which conserved quantities follow through the
characteristic polynomial. We obtain the representation from the
three-block construction of Jiang and Rauch-Wojciechowski
\cite[equation (2.7)]{JRW1991}, using the constant-diagonal extension
and the symmetric phase convention. The coordinate and phase changes
are given explicitly in Appendix~\ref{sec:JRW-dictionary}.

First write the equations using the mode matrices
$\boldsymbol u\in\C^{L\times M}$,
$\boldsymbol v\in\C^{M\times N}$,
$\boldsymbol w\in\C^{N\times L}$ and the constant phase matrices
\begin{equation}\label{eq:canonical-diagonals}
 \mathsf B_L=\diag(b_1^L,\ldots,b_L^L),\qquad
 \mathsf B_M=\diag(b_1^M,\ldots,b_M^M),\qquad
 \mathsf B_N=\diag(b_1^N,\ldots,b_N^N).
\end{equation}
Then \cref{eq:canonical-rates} is equivalently
\begin{equation}\label{eq:canonical-matrix-rates}
\begin{aligned}
 i\dot{\boldsymbol u}
 &=\mathsf B_L\boldsymbol u+\boldsymbol u\mathsf B_M
       +\sqrt2\bar{\boldsymbol w}^{\trans}\bar{\boldsymbol v}^{\trans},\\
 i\dot{\bar{\boldsymbol u}}^{\trans}
 &=-\mathsf B_M\bar{\boldsymbol u}^{\trans}-\bar{\boldsymbol u}^{\trans}\mathsf B_L
       -\sqrt2\boldsymbol v\boldsymbol w,\\
 i\dot{\boldsymbol v}
 &=\mathsf B_M\boldsymbol v+\boldsymbol v\mathsf B_N
       +\sqrt2\bar{\boldsymbol u}^{\trans}\bar{\boldsymbol w}^{\trans},\\
 i\dot{\bar{\boldsymbol v}}^{\trans}
 &=-\mathsf B_N\bar{\boldsymbol v}^{\trans}-\bar{\boldsymbol v}^{\trans}\mathsf B_M
       -\sqrt2\boldsymbol w\boldsymbol u,\\
 i\dot{\boldsymbol w}
 &=\mathsf B_N\boldsymbol w+\boldsymbol w\mathsf B_L
       +\sqrt2\bar{\boldsymbol v}^{\trans}\bar{\boldsymbol u}^{\trans},\\
 i\dot{\bar{\boldsymbol w}}^{\trans}
 &=-\mathsf B_L\bar{\boldsymbol w}^{\trans}-\bar{\boldsymbol w}^{\trans}\mathsf B_N
       -\sqrt2\boldsymbol u\boldsymbol v.
\end{aligned}
\end{equation}
These same blocks assemble into the $d\times d$ coefficient matrix
$\mathsf X(t)$, where $d=L+M+N$:
\begin{equation}\label{eq:canonical-X}
 \mathsf X=\begin{pmatrix}
 -2\mathsf B_L&\sqrt2\boldsymbol u&\bar{\boldsymbol w}^{\trans}\\
 \sqrt2\bar{\boldsymbol u}^{\trans}&-2\mathsf B_M&\boldsymbol v\\
 -\boldsymbol w&-\bar{\boldsymbol v}^{\trans}&\mathsf B_N
 \end{pmatrix}.
\end{equation}
The matrix pencil is the affine family
\begin{equation}\label{eq:canonical-pencil}
 \mathsf L(\nu,t)=\mathsf X(t)+\nu\mathsf D,\qquad
 \mathsf D=\diag(\Id_L,-\Id_M,0_N).
\end{equation}
Here $\nu$ is a constant spectral parameter, $\mathsf D$ is fixed,
and $\mathsf L(0,t)=\mathsf X(t)$. The time dependence of
$\mathsf X$ comes entirely from the mode matrices; its diagonal
blocks are constant. Suppressing $t$ in the notation, block
multiplication gives
\begin{equation}\label{eq:canonical-lax}
 i\dot{\mathsf L}=[\mathsf L,\mathsf T],\qquad
 \mathsf T(\nu)=\begin{pmatrix}
 -\mathsf B_L&0&\bar{\boldsymbol w}^{\trans}\\
 0&\mathsf B_M&-\boldsymbol v\\
 -\boldsymbol w&\bar{\boldsymbol v}^{\trans}&-\nu\Id_N
 \end{pmatrix}.
\end{equation}
Equivalently, $\dot{\mathsf L}=[\mathsf L,-i\mathsf T]$. Its off-diagonal blocks reproduce \cref{eq:canonical-matrix-rates},
while its diagonal blocks have zero derivatives.
The spectral determinant is
\begin{equation}\label{eq:canonical-spectral}
 F(\lambda,\nu)=\det(\lambda\Id_d-\mathsf L(\nu)).
\end{equation}
Its coefficients are polynomials in the modes and are conserved by
\cref{eq:canonical-lax}.

\subsection{The AKS framework and commuting integrals}
\label{sec:canonical-AKS}

To establish integrability, the conserved quantities must also commute
under the canonical bracket. The AKS framework explains this
commutativity by expressing the mode bracket in terms of the same
matrices $\mathsf X,\mathsf D$ used in
\cref{eq:canonical-X,eq:canonical-pencil}.
Write $X_{hi},D_{hi}$ for their entries. The canonical mode brackets give
\begin{equation}\label{eq:canonical-frozen}
 \{X_{hi},X_{jk}\}
 =-i(D_{hh}-D_{ii})\delta_{hk}\delta_{ij},
 \qquad1\leq h,i,j,k\leq d.
\end{equation}
Its coefficients depend only on the fixed matrix $\mathsf D$.
Fixing the three diagonal blocks of $\mathsf X$ leaves precisely the
off-diagonal canonical mode coordinates; these form a symplectic
affine leaf of the matrix bracket.

For scalar functions $f(\mathsf X),g(\mathsf X)$, define their
matrix gradients $P,Q$ by
\begin{equation}\label{eq:canonical-matrix-gradients}
 \mathrm df=\tr(P\,\mathrm d\mathsf X),\qquad
 \mathrm dg=\tr(Q\,\mathrm d\mathsf X).
\end{equation}
These gradients can depend on the mode values. Contracting
\cref{eq:canonical-frozen} gives
\begin{equation}\label{eq:canonical-gradient-bracket}
 \{f,g\}
 =i\tr(\mathsf D[P,Q]).
\end{equation}
This constant bracket and the shifted invariants of
$\mathsf X+\nu\mathsf D$ give the argument-shift realization of
the AKS construction \cite{MishchenkoFomenko1978,Bolsinov1992,
AdlerVanMoerbeke2018}.

The Hamiltonian belongs to this construction through the trace identity
\begin{align}\label{eq:canonical-trace-H}
 H={}&-\tfrac13\tr\mathsf X^3
      +3\tr\!\left[\mathsf B_L
        (\bar{\boldsymbol w}^{\trans}\boldsymbol w
                     -\boldsymbol u\bar{\boldsymbol u}^{\trans})\right]
        \nonumber\\
 &+3\tr\!\left[\mathsf B_M
        (\boldsymbol v\bar{\boldsymbol v}^{\trans}
                     -\bar{\boldsymbol u}^{\trans}\boldsymbol u)\right]
      -\tfrac83\tr\mathsf B_L^3
      -\tfrac83\tr\mathsf B_M^3
      +\tfrac13\tr\mathsf B_N^3.
\end{align}
The two quadratic trace terms are weighted combinations of the power
charges in \cref{eq:canonical-charges}; the traces of the fixed phase
matrices are constants. Thus $H$ combines a cubic spectral invariant
with commuting gauge generators.

\begin{proposition}\label{prop:canonical-commutation}
Every coefficient of $F$ in \cref{eq:canonical-spectral} commutes
with every other coefficient and with the charges in
\cref{eq:canonical-charges}.
\end{proposition}
\begin{proof}
Using the gradient definition \cref{eq:canonical-matrix-gradients},
let $P$ and $Q$ be the matrix gradients of conjugation-invariant polynomials
of $\mathsf X+\nu\mathsf D$ and $\mathsf X+\mu\mathsf D$,
respectively. They commute with their respective matrix arguments.
Cyclicity of the trace gives
\[
 \tr((\mathsf X+\nu\mathsf D)[P,Q])
 =\tr((\mathsf X+\mu\mathsf D)[P,Q])=0.
\]
Subtracting proves $\tr(\mathsf D[P,Q])=0$ for $\nu\ne\mu$;
polynomial continuation includes $\nu=\mu$. This proves spectral
commutativity under \cref{eq:canonical-frozen}. Each power charge
acts on $\mathsf X$ by constant diagonal conjugation, which also
preserves $F$.
\end{proof}

\subsection{The complete count for Liouville integrability}
\label{sec:canonical-independence}

The spectral invariants are complete when the phase coefficients are
distinct within each group. Equal phases allow additional continuous
changes of mode basis, whose conserved quantities supply any missing
integrals. We state the complete count here; Appendix~\ref{sec:integrability-proof}
gives its independence proof for arbitrary phase coefficients.

Write $s_h=D_{hh}$ and $e_h=X_{hh}$, $1\leq h\leq d$.
Partition each original group into maximal sets $B$ with a common
intercept $e_h$. We call these sets equal-phase blocks; their slopes
$s_h$ are also constant, with common value $s_B$.
Define the block moment matrix
\begin{equation}\label{eq:canonical-block-moment}
 (\mathsf J_B)_{ij}
 =\sum_{\substack{h=1\\s_h\ne s_B}}^d
       \frac{X_{ih}X_{hj}}{s_B-s_h},
 \qquad i,j\in B.
\end{equation}
For a constant $|B|\times|B|$ matrix $\mathsf Z$, let $\widehat{\mathsf Z}$ denote its
embedding by zeros outside $B$. Then
\[
 \{\mathsf X,\tr(\mathsf Z\mathsf J_B)\}=-i[\widehat{\mathsf Z},\mathsf X].
\]
Thus its Hamiltonian flow is conjugation by $e^{-i\epsilon\widehat{\mathsf Z}}$.
It preserves the Hamiltonian, since $\widehat{\mathsf Z}$ commutes with the fixed phase
matrices, so the entries of $\mathsf J_B$ are Noether integrals.
Different blocks act independently. Within each block the moment
entries have $-i$ times the matrix Lie--Poisson bracket.
These are the standard quadratic moment maps for rectangular canonical
matrix pairs \cite[Section~2]{Harnad1994}.

Fix an ordering of the indices in $B$, and let $\mathsf J_B^{[t]}$
denote its leading $t\times t$ principal submatrix. The functions
\begin{equation}\label{eq:canonical-moment-chain}
 I_{B,t,k}=\frac1k\tr\!\left((\mathsf J_B^{[t]})^k\right),
 \qquad 1\leq k\leq t\leq |B|,
\end{equation}
commute. For two nested principal blocks, the gradient of an
invariant of the larger block commutes with that block; the
Lie--Poisson pairing with the smaller-block gradient is therefore zero.
This is the Gelfand--Zeitlin construction
\cite{GuilleminSternberg1983}. These functions also commute with every
spectral coefficient, since block conjugation preserves $F$.
Let $\mathcal A$ be the polynomial algebra generated by these functions
and the coefficients of $F$. It is a commuting algebra of first
integrals of $H$.

\begin{theorem}\label{thm:canonical-independence}
For every positive $L,M,N$ and every fixed choice of phase coefficients,
$\mathcal A$ has generic differential rank $p=LM+MN+NL$.
Thus \cref{eq:canonical-H} is generically Liouville integrable.
A complete local set can include $H$ wherever $\mathrm dH\ne0$.
\end{theorem}
The proof in Appendix~\ref{sec:integrability-proof} proceeds by
increasing one matrix dimension at a time. At each insertion, the
spectral coefficients and principal moments together supply one new
independent integral per new canonical pair. This includes arbitrary
coincident phase coefficients. The same argument quantifies the
spectral deficit in \cref{eq:canonical-spectral-deficit}.

When all phases within each original group are distinct, the spectral
family alone is complete; an explicit selection by
line intersections and a spanning tree is given in
\cref{sec:distinct-phase-integrals}. Retaining the $d-1$ power charges
leaves
\begin{equation}\label{eq:canonical-genus-preview}
 G=p-d+1=LM+MN+NL-L-M-N+1
\end{equation}
spectral integrals on the regular locus used in
\cref{sec:spectral-geometry}.

Bolsinov's argument-shift completeness criterion
\cite[Theorem~1.1]{Bolsinov1992} includes the constant bracket
\cref{eq:canonical-frozen}. Its conditions are that
$\mathsf X+\nu\mathsf D$ is regular for every finite $\nu$ and that
the projected kernel has the required dimension; the latter becomes
distinctness of the phases within each original group.
A smooth spectral curve supplies the first condition, so completeness
in the distinct-phase case is a known consequence of this criterion.
\Cref{thm:canonical-independence} gives a uniform direct proof and
completes the count for repeated phases using the established
moment-map and Gelfand--Zeitlin constructions.
Jiang and Rauch-Wojciechowski prove completeness for row counts two
and three (our $L+M$), and leave four rows unresolved
\cite[Sections II.B--II.C]{JRW1991}. The insertion proof covers
arbitrary row and column counts and every fixed phase profile.

\begin{samepage}
\subsubsection{Commuting spectral flows and higher-order interactions}
\label{sec:spectral-flows}

The spectral invariants provide Hamiltonians as well as conserved
quantities. Choosing another invariant changes the evolution of the
same mode matrices while preserving their spectral curve. This is
the commuting-flow construction of the AKS framework
\cite{AdlerVanMoerbeke2018}, and applies to arbitrary positive
$L,M,N$.\par
\end{samepage}

For each integer $r\geq2$, define the trace Hamiltonian
\begin{equation}\label{eq:hierarchy-traces}
 H_r=\frac1r\tr\mathsf X^r.
\end{equation}
It is a spectral invariant at $\nu=0$ and is a polynomial of degree
at most $r$ in the modes. By
\cref{prop:canonical-commutation,eq:canonical-trace-H}, these
Hamiltonians commute with one another, with $H$ and $H_0$, and with
all the power charges and spectral coefficients. The complete
integral set in \cref{thm:canonical-independence} is therefore
preserved by every such flow.

For the $H_r$ evolution, the bracket
\cref{eq:canonical-frozen} gives an explicit Lax representation:
\begin{equation}\label{eq:hierarchy-lax}
 i\dot{\mathsf X}=[\mathsf D,\mathsf X^{r-1}],\qquad
 i\dot{\mathsf L}=[\mathsf L,\mathsf T_r],\qquad
 \mathsf T_r(\nu)=
       \frac{\mathsf X^{r-1}-\mathsf L(\nu)^{r-1}}{\nu}.
\end{equation}
The numerator vanishes at $\nu=0$, so $\mathsf T_r$ is polynomial
in $\nu$. Indeed, $[\mathsf L,\mathsf L^{r-1}]=0$ and
$[\mathsf X,\mathsf X^{r-1}]=0$ reduce its commutator to
$[\mathsf D,\mathsf X^{r-1}]$. The block-diagonal entries of
$\mathsf X$ remain fixed because $\mathsf D$ is scalar on each
of its three blocks.

The fourth member $H_4=\tr\mathsf X^4/4$ includes quartic
interactions, with lower-degree terms supplied by the fixed diagonal
entries. Appendix~\ref{sec:quartic-expansion} gives its expanded
quartic part.

For a constant $\alpha$, the combined Hamiltonian and time matrix are
\begin{equation}\label{eq:hierarchy-combination}
 H_\alpha=H+\alpha H_4,\qquad
 \mathsf T_\alpha=\mathsf T+\alpha\mathsf T_4.
\end{equation}
The same initial state and phase coefficients give the same
conserved polynomial $F(\lambda,\nu)$ under either evolution.
In a Jacobian or Prym linearization, if the geometric generators $i\partial_t$ have translation
vectors $\boldsymbol U_H,\boldsymbol U_4$, the divisor's Abel
coordinate evolves as
\begin{equation}\label{eq:hierarchy-direction}
 z(t)=z(0)-it(\boldsymbol U_H+\alpha\boldsymbol U_4).
\end{equation}
This is the spectral interpretation of choosing a different scalar
evolution within the commuting family
\cite{Beauville1990,Vanhaecke2001}.

The hyperelliptic realization in \cref{sec:BA-realization} selects
the scalar Schr\"odinger direction $z(t)=z_0-it e_g$ for $H_0$.
More generally, \cref{thm:BA-higher-flows} identifies the spectral
flows with multipoint BA times and constant reciprocal mode phases
for arbitrary group sizes. On each birational spectral level the
full family supplies every Abel direction.

